% Created 2020-02-18 Tue 20:38
% Intended LaTeX compiler: pdflatex

\documentclass[a4paper,12pt]{article}
\usepackage[margin=2cm]{geometry}

\usepackage{hyperref}
\usepackage[utf8]{inputenc}
\usepackage{fixltx2e}
\usepackage{graphicx}
\usepackage{longtable}
\usepackage{float}
\usepackage{wrapfig}
\usepackage{rotating}
\usepackage[normalem]{ulem}
\usepackage{amsmath}
\usepackage{textcomp}
\usepackage{marvosym}
\usepackage{wasysym}
\usepackage{multicol}
\usepackage{amssymb}
\tolerance=1000
\usepackage{listings}
\usepackage{xeCJK}
\setCJKmainfont{SimSun}
\usepackage{color}
\usepackage{listings}
\RequirePackage{fancyvrb}
\DefineVerbatimEnvironment{verbatim}{Verbatim}{fontsize=\scriptsize}
\lstset{basicstyle=\scriptsize\ttfamily}
\lstset{basicstyle=\tiny\ttfamily}
\usepackage{placeins}
\vspace{-0.2cm}
\usepackage{xcolor}
\definecolor{dkgreen}{rgb}{0,0.6,0}
\definecolor{dred}{rgb}{0.545,0,0}
\definecolor{dblue}{rgb}{0,0,0.545}
\definecolor{lgrey}{rgb}{0.9,0.9,0.9}
\definecolor{gray}{rgb}{0.4,0.4,0.4}
\definecolor{darkblue}{rgb}{0.0,0.0,0.6}
\definecolor{bubbles}{rgb}{0.91, 1.0, 1.0}
\lstdefinelanguage{python}{
backgroundcolor=\color{bubbles},
basicstyle=\footnotesize \ttfamily \color{black} \bfseries,
breakatwhitespace=false,
breaklines=true,
captionpos=b,
commentstyle=\color{dred},
deletekeywords={...},
escapeinside={\%*}{*)},
frame=single,
language=C++,
keywordstyle=\color{purple},
morekeywords={BRIEFDescriptorConfig,string,TiXmlNode,DetectorDescriptorConfigContainer,istringstream,cerr,exit},
identifierstyle=\color{black},
stringstyle=\color{blue},
rulecolor=\color{black},
showspaces=false,
showstringspaces=false,
showtabs=false,
stepnumber=1,
tabsize=5,
title=\lstname,
}
\usepackage{enumitem}
\usepackage[utf8]{inputenc}
\author{yen yung chin}
\date{\today}
\title{Python programming}
\begin{document}

\maketitle
\tableofcontents



\section{Python程式設計}
\label{sec:orgfb5ffd5}
\subsection{變數、資料型態、輸出輸入}
\label{sec:org88bb171}
\subsubsection{變數}
\label{sec:orgc596e1c}
\begin{enumerate}
\item 命名規則
\label{sec:org315aad5}
\begin{itemize}itemsep=1pt
\item 由英文、數字、底線、中文(不建議)組成\\
\item 不得以數字開頭\\
\item 不能與Python內建的保留字相同\\
\end{itemize}
\item practice
\label{sec:org66743b3}
\begin{itemize}itemsep=2pt
\item abc\_123\\
\item 3pigs\\
\item Happy New Year\\
\item Class\\
\item Good!\\
\end{itemize}
\end{enumerate}

\subsubsection{變數的指派(assign)}
\label{sec:orge4a9153}
\begin{itemize}
\item Python變數不需宣告，依指派值自動設定資料型態。\\
\end{itemize}

\begin{enumerate}
\item 指派語法：

\label{sec:orgc24164c}
\begin{itemize}
\item 變數名稱 = 指派值\\
\item 變數不再使用時，可用del指令將其刪除，以節省記憶體。\\
\end{itemize}

\item 範例：
\label{sec:orgd93b46c}
\lstset{breaklines=true,language=Python,label= ,caption= ,captionpos=b,firstnumber=1,numbers=left}
\begin{lstlisting}
a = 5
b = 3.14
c = 'TNFSH'
a = b = c = 10
name, number = 'TNFSH', 35
\end{lstlisting}
\end{enumerate}

\subsubsection{資料型態}
\label{sec:org5ae6e7e}
\begin{enumerate}
\item 常用類型
\label{sec:org06799f4}
\begin{itemize}
\item 整數: int\\
\item 浮點數: float\\
\item 布林值(True / False): bool, T 與 F要大寫\\
\item 字串: str, 以'或"含括,若輸出字串要包含引號，則以另一種引號含括該字串。如:\\
\end{itemize}
\lstset{breaklines=true,language=Python,label= ,caption= ,captionpos=b,firstnumber=1,numbers=left}
\begin{lstlisting}
good_words = " 請常說'請'、'謝謝'、'對不起' "
print(good_words)
\end{lstlisting}

\begin{verbatim}
 請常說'請'、'謝謝'、'對不起' 
\end{verbatim}

\item 型態間的轉換
\label{sec:org9a7d4eb}

\begin{itemize}
\item 自動轉換\\
\end{itemize}
\lstset{breaklines=true,language=Python,label= ,caption= ,captionpos=b,numbers=none}
\begin{lstlisting}
score = 60 + 3.5		# 自動轉換為浮點數，結果為63.5	
\end{lstlisting}

\begin{itemize}
\item 強制轉換\\
\end{itemize}
\lstset{breaklines=true,language=Python,label= ,caption= ,captionpos=b,numbers=none}
\begin{lstlisting}
score = int(30.22)			# 將括弧內的資料轉換為整數
score = float(score)		# 將括弧內的資料轉換為浮點數
test = str(score)				# 將括弧內的資料轉換為字串
\end{lstlisting}
\end{enumerate}

\subsubsection{輸出: print}
\label{sec:orge21230c}
\begin{enumerate}
\item 語法
\label{sec:orgdbfd770}
\begin{itemize}
\item print( 項目1, [ 項目2, … , sep = 分隔字元 , end = 結束字元 ] )\\
\item sep (分隔字元) 預設為空白字元\\
\item end (結束字元) 預設為換行字元\n \\
\end{itemize}
\item 範例
\label{sec:org43c9a9b}
\lstset{breaklines=true,language=Python,label= ,caption= ,captionpos=b,firstnumber=1,numbers=left}
\begin{lstlisting}
a, b = 5, 10
print(a, b)
print(a, b, sep =',')
print(a, end='\t')
print(b)
print(a, end='')
print(b)
\end{lstlisting}

\begin{verbatim}
5 10
5,10
5	10
510
\end{verbatim}
\item escape char
\label{sec:org196faff}
\begin{center}
\begin{tabular}{ll}
\hline
跳脫字元 & 說明\\
\hline
$\backslash$\ & 印出 $\backslash$\\
$\backslash$' & 印出 '\\
\n & 換行字元\\
\t & 水平跳格字元(Tab)\\
\b & 倒退一格字元(Backspace)\\
\hline
\end{tabular}
\end{center}

\item 格式化輸出
\label{sec:orgacb6e1d}
\begin{enumerate}
\item 語法
\label{sec:org73436d5}
\begin{itemize}
\item 格式化輸出語法：print( 字串  \%(參數) )
\\
\item PS: 字串裡 \%s 代表字串、\%d 代表整數、\%f 代表浮點數\\
\end{itemize}
\item 範例
\label{sec:org9fe253a}
\lstset{breaklines=true,language=Python,label= ,caption= ,captionpos=b,firstnumber=1,numbers=left}
\begin{lstlisting}
name = '台北101'
height = 508
fee = 18.24
print('%s的高度為%d公尺，參觀門票為%3.2f美金.' %(name, height, fee))

\end{lstlisting}

\begin{verbatim}
台北101的高度為508公尺，參觀門票為18.24美金.
\end{verbatim}
\end{enumerate}

\item format格式化輸出
\label{sec:orgf711221}
\begin{enumerate}
\item 語法
\label{sec:org21f2e89}
\begin{itemize}
\item format格式化輸出語法：print(字串.format (參數) )
\\
\item PS: 字串裡以\{0\}、\{1\}、\ldots{} 來對應參數列裡的變數\\
\end{itemize}
\item 範例
\label{sec:orgc8c06d3}
\lstset{breaklines=true,language=Python,label= ,caption= ,captionpos=b,firstnumber=1,numbers=left}
\begin{lstlisting}
name = 'Taipei101'
height = 508.35
print('{0}的高度為{1}公尺'.format (name, height))
print('{0:15s}的高度為{1:10.1f}公尺'.format (name, height))
\end{lstlisting}

\begin{verbatim}
Taipei101的高度為508.35公尺
Taipei101      的高度為     508.4公尺
\end{verbatim}
\end{enumerate}
\end{enumerate}

\subsubsection{輸入}
\label{sec:org4801b6a}
\begin{enumerate}
\item 語法
\label{sec:orged74e33}
\begin{itemize}
\item 輸入語法：變數 = input( [提示字串] )
\\
\item PS: 經由input( )函式讀入的資料，其資料型態皆為字串\\
\end{itemize}

\item 範例
\label{sec:orgd83b110}
\lstset{breaklines=true,language=Python,label= ,caption= ,captionpos=b,firstnumber=1,numbers=left}
\begin{lstlisting}
a = input('輸入國文成績:')
b = input('輸入數學成績:')
c = input()
print('三科成績分別為%5s %5s %5s' %(a, b, c))
\end{lstlisting}

\item 輸入搭配型別轉換
\label{sec:org61b12f5}
\begin{enumerate}
\item 範例
\label{sec:orgc112a56}
\lstset{breaklines=true,language=Python,label= ,caption= ,captionpos=b,firstnumber=1,numbers=left}
\begin{lstlisting}
a = int(input('輸入國文成績:'))
b = int(input('輸入數學成績:'))
c = int(input())
print('總分為%5d' %(a+b+c))
\end{lstlisting}
\end{enumerate}

\item input注意事項
\label{sec:orgeb00edd}
\begin{itemize}
\item input 每筆輸入讀至換行符號為止\\
\item 若測試資料輸入以空白間隔，則讀入整行字串後，
再以字串分割處理\\
\end{itemize}
\lstset{breaklines=true,language=Python,label= ,caption= ,captionpos=b,firstnumber=1,numbers=left}
\begin{lstlisting}
inp = input()
inpstr = inp.split()
# 依此類推逐一取得輸入值
a = int(inpstr[0])
b = int(inpstr[1])
c = int(inpstr[2])
\end{lstlisting}
\end{enumerate}

\subsubsection{註解}
\label{sec:org0267d06}
\begin{enumerate}
\item 語法
\label{sec:orgff1de68}
\begin{itemize}
\item 單行註解：以 \# 開頭\\
\item 多行註解：前後以 ''' 或 """ 含括\\
\end{itemize}

\item 範例
\label{sec:orgad58e0d}
\lstset{breaklines=true,language=Python,label= ,caption= ,captionpos=b,firstnumber=1,numbers=left}
\begin{lstlisting}
## 這是單行註解
int a = 3
'''
這是多行註解
LALALA
'''
print(a)
"""
這也是
"""
\end{lstlisting}
\end{enumerate}

\subsubsection{實作練習}
\label{sec:org89c46e6}
\begin{itemize}
\item \href{http://skyoj.tnfsh.tn.edu.tw/sky/index.php/problem/view/161/}{skyoj \#161華氏溫度轉攝氏溫度}\\
\end{itemize}
\subsection{運算元與運算式}
\label{sec:org8ef5b3a}
\subsubsection{算術運算}
\label{sec:orgd1df183}
\begin{itemize}
\item +: 加\\
\item -: 減\\
\item *: 乘\\
\item /: 除\\
\item \%: 取餘數\\
\item //: 求商\\
\item **: 指數\\
\end{itemize}

\subsubsection{範例}
\label{sec:orgd1a297a}
\lstset{breaklines=true,language=Python,label= ,caption= ,captionpos=b,firstnumber=1,numbers=left}
\begin{lstlisting}
# python code for arithematic opearation
print(5+3)
print(5-3)
print(5*3)
print(5/3)
print(5%3)
print(5//3)
print(5**3)
\end{lstlisting}

\begin{verbatim}
Hello world
Nice to meet you, Letranger
\end{verbatim}

\subsubsection{input 指令}
\label{sec:org10aafd3}
\subsubsection{if else 語法}
\label{sec:org8dd342d}
\subsubsection{四則運算}
\label{sec:orgcf04b7e}
\begin{itemize}
\item // 整除\\
\item \% 取餘數\\
\end{itemize}
\lstset{breaklines=true,language=Python,label=計算機,caption= ,captionpos=b,firstnumber=1,numbers=left}
\begin{lstlisting}
a = int(input("input a: "))
op = input("input op: ")
b = int(input("input b: "))

if op == '+':
    ans = a + b;
elif op == '-':
    ans = a - b;
elif op == '*':
    ans = a * b;
elif op == '/':
    ans = a / b;
elif op == '%':
    ans = a % b;

print(str(a) + op + str(b) + "=" + str(ans))
\end{lstlisting}

\subsubsection{由a加到b}
\label{sec:org96acf98}
\lstset{breaklines=true,language=Python,label=計算機,caption= ,captionpos=b,firstnumber=1,numbers=left}
\begin{lstlisting}
a = int(input("input a: "))
b = int(input("input b: "))
# while solution
sum = 0
while a <= b:
    sum = sum + a
    a = a + 1
print("Result :" + str(sum))
# for solution
sum = 0
for i in range(a, b):
    sum = sum + i
print("Result:" + str(sum))
\end{lstlisting}

\subsubsection{由a乘到b}
\label{sec:org9a75bfe}
\lstset{breaklines=true,language=Python,label=計算機,caption= ,captionpos=b,firstnumber=1,numbers=left}
\begin{lstlisting}
a = int(input("input a: "))
b = int(input("input b: "))
# for solution
x = 1
for i in range(a, b + 1):
    x = x * i
print(x)
# while solution
p = 1
while a <= b:
    p = p * a
    a = a + 1

print("Result:" + str(p))
\end{lstlisting}

\subsubsection{list}
\label{sec:org06ffc3e}
\lstset{breaklines=true,language=sh,label= ,caption= ,captionpos=b,numbers=none}
\begin{lstlisting}
date
\end{lstlisting}
\lstset{breaklines=true,language=Python,label= ,caption= ,captionpos=b,firstnumber=1,numbers=left}
\begin{lstlisting}

import urllib.request as req
url = "https://huodalife.pixnet.net/blog/"
# 幫request加上一個header
request = req.Request(url, headers = {
    "User-Agent":"Mozilla/5.0 (Macintosh; Intel Mac OS X 10.14; rv:72.0) Gecko/20100101 Firefox/72.0"
})

with req.urlopen(request) as response:
    data = response.read().decode("utf-8")

import bs4
root = bs4.BeautifulSoup(data,"html.parser")

titles = root.find_all("li", class_="title")
pubs = root.find_all("li", class_="publish")
for pub, title in zip(pubs, titles):
    print(pub.span.string,title.a.string)
\end{lstlisting}


\subsection{0725}
\label{sec:orgc613a68}
\subsubsection{swap a, b}
\label{sec:org64df84e}
\begin{itemize}
\item a, b = b, a\\
\end{itemize}
\subsubsection{gcd}
\label{sec:org3ba7df7}
\lstset{breaklines=true,language=Python,label= ,caption= ,captionpos=b,firstnumber=1,numbers=left}
\begin{lstlisting}
a = input()
b = input()
while(a % b != 0):
    b = a
    a = a % b
print(b)
\end{lstlisting}

\subsubsection{python 14 list}
\label{sec:org5252041}
\lstset{breaklines=true,language=Python,label= ,caption= ,captionpos=b,firstnumber=1,numbers=left}
\begin{lstlisting}
import numpy as np
students = []
grades = []
for i in range(6):
    name = input()
    grades.append(name)
    for j in range(3):


\end{lstlisting}

\subsection{0802 CNN}
\label{sec:orge312555}
\url{https://reurl.cc/Y2gKo}: 各種常用的library, function\\

\section{以python下載網路資料}
\label{sec:org46394e1}
\begin{itemize}
\item import urllib.request as req\\
\end{itemize}

\section{網路爬蟲}
\label{sec:org32a7038}
\subsection{parse HTML}
\label{sec:org758b9b4}
\begin{itemize}
\item what is HTML\\
\item 精神：將程式模仿成一般user\\
\end{itemize}
\subsubsection{所需套件}
\label{sec:org536b198}
\begin{itemize}
\item BeautifulSoup\\
\end{itemize}
pip install beautifulsoup4\\
\subsection{範例: 以程式抓取PTT電影板header}
\label{sec:org91a8952}
\subsubsection{原則}
\label{sec:orgf5955de}
\begin{enumerate}
\item 要抓網路資料，就要先仔細觀察網頁內容\\
\item 要儘量欺騙網站自己是browser\\
\end{enumerate}
\subsubsection{實作: 取得HTML資料}
\label{sec:org83b4602}
\begin{enumerate}
\item 第一版
\label{sec:org728b532}
\lstset{breaklines=true,language=Python,label= ,caption= ,captionpos=b,firstnumber=1,numbers=left}
\begin{lstlisting}
# 抓取PTT電影版header
import urllib.request as req
url = "https://www.ptt.cc/bbs/Stock/index.html"

with req.urlopen(url) as response:
    data = response.read().decode("utf-8")
print(data)
\end{lstlisting}

\begin{verbatim}
urllib.error.HTTPError: HTTP Error 403: Forbidden
\end{verbatim}

被拒絕原因：這隻程式的行為不像一般使用者，被網站伺服器拒絕。403 Forbidden 是HTTP協議中的一個HTTP狀態碼（Status Code）。可以簡單的理解為沒有權限訪問此站，服務器收到請求但拒絕提供服務\footnote{\href{https://zh.wikipedia.org/wiki/HTTP\_403}{維基百科: HTTP 403} \\}。\\

\item 第二版:加入header
\label{sec:orgc3fa9dd}
\lstset{breaklines=true,language=Python,label= ,caption= ,captionpos=b,firstnumber=1,numbers=left}
\begin{lstlisting}
import urllib.request as req
url = "https://www.ptt.cc/bbs/Stock/index.html"
# 幫request加上一個header
request = req.Request(url, headers = {
    "User-Agent":"Mozilla/5.0 (Macintosh; Intel Mac OS X 10.14; rv:72.0) Gecko/20100101 Firefox/72.0"
})

with req.urlopen(request) as response:
    data = response.read().decode("utf-8")
print('所取得的資料型態：',type(data))
print(data[700:1000])
\end{lstlisting}

\begin{verbatim}
所取得的資料型態： <class 'str'>
er">
	<div id="topbar" class="bbs-content">
		<a id="logo" href="/bbs/">批踢踢實業坊</a>
		<span>&rsaquo;</span>
		<a class="board" href="/bbs/Stock/index.html"><span class="board-label">看板 </span>Stock</a>
		<a class="right small" href="/about.html">關於我們</a>
		<a class="right small" href="/contact.html">
\end{verbatim}
\end{enumerate}

\subsubsection{實作：解析HTML資料內容}
\label{sec:org7de1d18}
\begin{enumerate}
\item 安裝解析HTML所需套件
\label{sec:org64c0039}
\lstset{breaklines=true,language=shell,label= ,caption= ,captionpos=b,numbers=none}
\begin{lstlisting}
pip install beautifulsoup4
\end{lstlisting}

\item 第三版: 解析HTML
\label{sec:orgfac7f19}
\lstset{breaklines=true,language=Python,label= ,caption= ,captionpos=b,firstnumber=1,numbers=left}
\begin{lstlisting}
import urllib.request as req
url = "https://www.ptt.cc/bbs/Stock/index.html"
# 幫request加上一個header
request = req.Request(url, headers = {
    "User-Agent":"Mozilla/5.0 (Macintosh; Intel Mac OS X 10.14; rv:72.0) Gecko/20100101 Firefox/72.0"
})

with req.urlopen(request) as response:
    data = response.read().decode("utf-8")

import bs4
root = bs4.BeautifulSoup(data,"html.parser")
print(root.title)
titles = root.find("div", class_="title")
print(titles)
print(titles.a.string)
titles = root.find_all("div", class_="title")
print(titles)
for title in titles:
    if title.a != None:
        print(title.a.string)
\end{lstlisting}

\begin{verbatim}
<title>看板 Stock 文章列表 - 批踢踢實業坊</title>
<div class="title">
<a href="/bbs/Stock/M.1581985895.A.742.html">Re: [新聞] 蝗蟲抵達邊境 中國人自信：威脅不大</a>
</div>
Re: [新聞] 蝗蟲抵達邊境 中國人自信：威脅不大
[<div class="title">
<a href="/bbs/Stock/M.1581985895.A.742.html">Re: [新聞] 蝗蟲抵達邊境 中國人自信：威脅不大</a>
</div>, <div class="title">
<a href="/bbs/Stock/M.1581986080.A.0B9.html">[新聞] 武漢肺炎衝擊 蘋果率先示警營收無法達標</a>
</div>, <div class="title">
<a href="/bbs/Stock/M.1581987000.A.EE7.html">[新聞] 超預期 國巨MLCC漲價五成</a>
</div>, <div class="title">
<a href="/bbs/Stock/M.1581987018.A.192.html">[新聞] 武漢肺炎》上海美國商會：1/3企業若工廠</a>
</div>, <div class="title">
<a href="/bbs/Stock/M.1581987846.A.95F.html">[新聞] 中芯國際14nm第二代FinFET批量須等到2020</a>
</div>, <div class="title">
<a href="/bbs/Stock/M.1581989947.A.392.html">Re: [請益] 有人買股票而完全捨棄定存嗎?</a>
</div>, <div class="title">
<a href="/bbs/Stock/M.1581990753.A.EB3.html">[新聞] 斷絕華為貨源！美擬祭禁令限制中企取得</a>
</div>, <div class="title">
<a href="/bbs/Stock/M.1581994926.A.A50.html">[標的] 6531愛普</a>
</div>, <div class="title">
<a href="/bbs/Stock/M.1581996309.A.C55.html">[請益] 是不是該立法更限制外資匯出資金？</a>
</div>, <div class="title">
<a href="/bbs/Stock/M.1581999672.A.14E.html">Re: [請益] 存股標的選擇</a>
</div>, <div class="title">
<a href="/bbs/Stock/M.1582000361.A.6A2.html">[標的] 1201 味全 多</a>
</div>, <div class="title">
<a href="/bbs/Stock/M.1582000701.A.2AB.html">Re: [新聞] 蝗蟲抵達邊境 中國人自信：威脅不大</a>
</div>, <div class="title">
<a href="/bbs/Stock/M.1582001420.A.3B8.html">[新聞] 用電大戶買綠電 電證分離更活絡</a>
</div>, <div class="title">
<a href="/bbs/Stock/M.1422199105.A.84E.html">[公告] 精華區導覽Q&amp;A</a>
</div>, <div class="title">
<a href="/bbs/Stock/M.1541252211.A.8BE.html">[公告] Stock 板規V2.6             (2019/06/20)</a>
</div>, <div class="title">
<a href="/bbs/Stock/M.1579935755.A.395.html">[公告] 關於武漢肺炎新聞規範</a>
</div>, <div class="title">
<a href="/bbs/Stock/M.1581919202.A.CF2.html">[閒聊] 2020/02/17 盤後閒聊</a>
</div>, <div class="title">
<a href="/bbs/Stock/M.1581985803.A.F3E.html">[閒聊] 2020/02/18 盤中閒聊</a>
</div>]
Re: [新聞] 蝗蟲抵達邊境 中國人自信：威脅不大
[新聞] 武漢肺炎衝擊 蘋果率先示警營收無法達標
[新聞] 超預期 國巨MLCC漲價五成
[新聞] 武漢肺炎》上海美國商會：1/3企業若工廠
[新聞] 中芯國際14nm第二代FinFET批量須等到2020
Re: [請益] 有人買股票而完全捨棄定存嗎?
[新聞] 斷絕華為貨源！美擬祭禁令限制中企取得
[標的] 6531愛普
[請益] 是不是該立法更限制外資匯出資金？
Re: [請益] 存股標的選擇
[標的] 1201 味全 多
Re: [新聞] 蝗蟲抵達邊境 中國人自信：威脅不大
[新聞] 用電大戶買綠電 電證分離更活絡
[公告] 精華區導覽Q&A
[公告] Stock 板規V2.6             (2019/06/20)
[公告] 關於武漢肺炎新聞規範
[閒聊] 2020/02/17 盤後閒聊
[閒聊] 2020/02/18 盤中閒聊
\end{verbatim}
\end{enumerate}

\subsection{參考資料}
\label{sec:orgf064e47}
\begin{itemize}
\item \url{https://www.youtube.com/watch?v=9Z9xKWfNo7k}\\
\end{itemize}


\section{如何把python程式publish成網站應用程式}
\label{sec:org2927c7e}
\end{document}
